\section{Discussion}\label{sec:discussion}

\subsection{What is established and what is hypothetical}
Three things are established.
\begin{enumerate}
\item The refined crosscap invariant $g = f - \tfrac12 f^{\cov}$, in the uniform normalization and hatted with $M(q^{-1})$, is an integer Laurent polynomial in every case examined (unknot, Hopf link, torus knots in the fundamental representation, the trefoil with two to four boxes, $T(2,5)$ with two and three boxes, $T(2,7)$ and $T(2,9)$ with three boxes, algebraic cables in the fundamental representation and with two boxes; for four boxes $2(1+t)g$), and reduces at $t = q$ to the results of \cite{Marino09} for the fundamental representation and the two-box trefoil [computation].
\item The all-$d$ structure $g = \tfrac12 E/(1+t) + C$, the appearance of $E$ only in even sectors, and the boundary factor $z_q^{n-1}(z_q^+)^n$ of $E$ are derived under Hypothesis \ref{hyp:classes} [derivation] and agree with computation [computation].
\item The brane--brane sector of the cover is the invariant of \cite{KN17} itself, and the brane--image sector is identified as a space of states [computation].
\end{enumerate}
What remains hypothetical: that the cover is given by A-type refined invariants (Hypothesis \ref{hyp:hg}; supported by the integrality of $g$ after subtraction and by the uniqueness of the pole cancellation), that the uniform normalization is framing zero (Hypothesis \ref{hyp:framing}), and the classification of one-particle states (Hypothesis \ref{hyp:classes}).
As conjectures we put forward the criterion in the fundamental representation (Conjecture \ref{conj:criterion}; checked for $T(2,2m+1)$, $m \le 5$, and restricted to two-strand torus knots since it fails for $T(3,4)$ and for algebraic cables) and the structure with the universal factor (Conjecture \ref{conj:structure}).

\subsection{Relation to \texorpdfstring{\cite{KN17}}{[KN17]}}\label{sec:kn}
Positivity here means the criterion of \S\ref{sec:criterion}: after the $G$-expansion, the signs align within each $\lambda^\beta$ block (\S\ref{sec:positivity}).
The brane--brane sector of the cover (the oriented part) agrees coefficient by coefficient with the $\U(N_A)$ refined LMOV invariants of \cite{KN17} (\S\ref{sec:split}) and therefore inherits their $G$-expansion and positivity (Appendix \ref{app:cover}).
The brane--image sector, after removal of the boundary factors, is aligned except for the $-S\otimes S$ in the switching block $k = 2m$, and that mixing is understood as the subtraction of pairs of states (\S\ref{sec:split}).
What breaks is the multi-box crosscap part $g$, the sector of fixed points of $\tilde I$, which carries the same universal factor $(\bt^2 - 1)$.
The evidence covers the trefoil and $T(2,5)$ with two boxes ($T(2,7)$ for brane--image); the three-box invariants have not been split into brane and image.
Thus the positivity problem for refined Kauffman invariants reduces (conjecturally) to adding to the framework of \cite{KN17} the parity structure of the orientifold: the sign of $\tilde I$ and the multiplet basis of the frozen states.

\subsection{Open problems}
\begin{enumerate}
\item \textbf{Multiplets of frozen states.} The multi-box $c = 1$ part is not $\sigma$-symmetric and cannot be expanded in $G$; the factors $T_\pm$ of $G = T_+T_-$ do not help (remark in \S\ref{sec:index}).
By Result \ref{cr:multibox-sigma}, the unrefined component of the asymmetry is removed by the parity rule for $c$, and the remaining refined $\sigma$-odd part is divisible by $(\bt^2-1)(\ql - \ql^{-1})$. The $\sigma$-even/odd split is half-integral and the signs do not align for $T(2,5)$, so this split does not suffice for positivity. Increasing the knot to $T(2,9)$ with three boxes and the number of boxes to four for the trefoil preserves integrality, the position of the $(1+t)$ poles and the universal factor, while the fraction of aligned blocks decreases (Result \ref{cr:signs}).
The $G$-expansion originates in the structure of oriented surfaces (the Lefschetz $\mathfrak{sl}_2$ on the Hodge decomposition of $H^1$), and there is no reason to expect it for frozen states wrapped on non-orientable surfaces (\S\ref{sec:index}). Determining the refined handle and crosscap factors of non-orientable surfaces from geometry is the first task.
\item \textbf{Separation of $E$ and $C$.} It is not determined by the index (Proposition \ref{prop:EC}); it is undetermined for the Hopf link $(4,4)$ and for knots with two boxes (for brane 2, image 2 of the trefoil it is fixed up to $\lambda^{-12}$ and undetermined from $\lambda^{-14}$ on, Appendix \ref{app:n2}).
The action on boundary zero modes is fixed (Proposition \ref{prop:boundary}). What remains is the moduli part: fixing the sign of $\tilde I = I_X\otimes e^{i\pi J_2}$ requires splitting the $t$-charge of a state into $J_2$ and $S_R$, which the index does not provide. That the cohomological Hall module of \cite{Young16} is not free over the CoHA is the algebraic form of this ambiguity (for the loop quiver it is determined, consistent with the uniqueness of $E$ for the framed unknot).
\item \textbf{Meaning of the universal factor.} $(\ba^2 - \bq^2)$ is the $\SO(2)$ point and $(\bt^2 - 1)$ is explained by $t = q$ and parity. Why the direction of the shift $\lambda \to \tlam$ is unique, and its relation to the tripling of the closed string K\"ahler variable, remain unexplained.
\item \textbf{Other knots.} For non-torus knots the $S, T$ matrices of refined Chern--Simons theory are not available, but for iterated torus knots there is the DAHA of \cite{CD14}, whose $D_N$ values we adopted as the definition for two algebraic cables (\S\ref{sec:cable}). In the fundamental representation the criterion fails there; what survives is integrality and the universal factor. Why the $c = 1$ part is $\sigma$-symmetric only for two-strand torus knots is not understood \open. Non-algebraic cables were not treated. The two- and three-box invariants of the $(2,15)$-cable are out of reach of the present computation (the $D_N$ basis at $N \ge 13$ is estimated to exceed 60\,GB).
At fixed $N$, the identifications $\SO(4) = \mathrm{SU}(2)^2$, $\SO(6) = \mathrm{SU}(4)$ allow A-type input (e.g., the rectangular formulas of \cite{KNTZ19}); for seven double twist knots we checked the integrality of the hatted $\SO$ invariants at $N = 2$ ($|R| \le 3$) and the uniform normalization.
\item \textbf{Type $B$ and higher colors of $\Sp$.} $\Sp(2N)$ was treated in \S\ref{sec:sp} (fundamental representation and the two-box trefoil). The two-box $T(2,5)$ needs $N \le 13$ for the interpolation, too costly for the present implementation. The closed string of type $B$ \cite{AM22} has only $c = 1$ and the denominator $(t^{n/2} + t^{-n/2})$; how type $B$ fits the classification of Proposition \ref{prop:alld} has not been examined.
\item \textbf{Hilbert scheme models.} Implementing the cell decomposition of \cite{ORS12} for $\Gamma = \langle 2, 3\rangle$ reproduces the HOMFLY homology of $T(2,3)$ (the unreduced superpolynomial of \cite{DGR05}) up to degree 16 in $q^{2l}$ and fixes the dictionary between $(q, t, a)$ degrees and the (length, weight, number of generators) of cells. On the other hand, the index of the fixed sector $P = (1-t)x$ factorizes as $-(\ba^2 - 1)(\ba^2 - \bq^2)(\bq^2\bt^2 - 1)(\ba^2\bq^2\bt^2 - 1) N_8/(\ba^{12}\bq^4\bt^{10})$ ($N_8$ has eight terms and the same boundary factor $(\lambda^2-1)(q-1)$ as the framed unknot) and is not the $\psi^2$ of the fundamental representation of the cover. To fix $E$ from geometry one needs a cell model of the brane--image composite sector together with a real structure on it (computation and conjecture).
\end{enumerate}

\subsection{Outlook}
The definition and the structures of this paper reach a closed-form criterion for the fundamental representation of two-strand torus knots, and integrality and the universal factor for higher representations and for algebraic cables; the range of the definition has been extended to algebraic cables through the DAHA of \cite{CD14} (\S\ref{sec:cable}).
The next steps are to determine the multiplet structure of frozen states from the geometry of non-orientable surfaces, and to understand why the criterion in the fundamental representation is specific to two-strand torus knots. Once the former is settled, the multi-box sign problem of \S\ref{sec:positivity} can be formulated as a multi-box version of the criterion.
